%%%%%%%%%%%%%%%%%%%%%%%%%%%%%%%%%%%%%%%%%%%%%%%%
% COPYRIGHT: (C) 2012-2026 FAU FabLab and others
% Lizenz: CC-BY-SA 3.0 (wie alle anderen Einweisungen des FAU FabLab)
%%%%%%%%%%%%%%%%%%%%%%%%%%%%%%%%%%%%%%%%%%%%%%%%


\PassOptionsToPackage{english}{babel} % zusätzlich zu ngerman aus fablab-document
\newcommand{\basedir}{fablab-document}
\documentclass[13pt]{\basedir/fablab-document}

\usepackage{wrapfig} %Textumlauf um Bilder
\date{\today}
\author{kontakt@fablab.fau.de}
\title{General Workshop Rules}
\fancyfoot[C]{Page \thepage~of~\pageref*{LastPage}}
\linespread{1} % Zeilenabstand

\usepackage{parskip} % Abstände zwischen Absätzen / Listenelementen
 \setlength{\parskip}{0.7\parskip}
% \usepackage[compact]{titlesec} % Abstände bei Überschriften
% \titlespacing{\section}{0pt}{*0.5}{*0}
% \titlespacing{\subsection}{0pt}{*0.3}{*0}
%\titlespacing{\subsubsection}{0pt}{*0}{*0}

%\fancyfoot[L]{}
%\fancyfoot[C]{}
%\fancyfoot[R]{Version 2, April 2012}

\setcounter{secnumdepth}{0}

% solche Wörter werden nicht getrennt!


% Neuer Befehl \subscript (Text tiefgestellt) von http://anthony.liekens.net/index.php/LaTeX/SubscriptAndSuperscriptInTextMode
%\newcommand{\subscript}[1]{\ensuremath{_{\textrm{\small{#1}}}}}

\begin{document}
\selectlanguage{english}

\maketitle

\vbox{\vspace{1cm}}


\section{Code of conduct}
\begin{itemize}
  \item Independent work in the FabLab rooms usually requires an instruction
  \item \textbf{No smoking} anywhere in the FabLab.
  \item Treat machines and tools with care! Inspect tools and equipment visually before use! Hand in anything defective to a supervisor (\enquote{Betreuer}) right away. Defective items will be repaired by a qualified person or thrown away.
  \item \textbf{Keep the place tidy!} Throw away dirt and leftovers.
  \item Only supervisors may store things in the FabLab. Unlabeled items lying around will be disposed of or reused as donations.
  \item \textbf{Work carefully}, ask if anything is unclear
  \item In case of accidents
  \begin{itemize}
  	\item Notify a supervisor (don't worry, we won't be angry with you if something happens!)
  	\item In case of serious injury or danger to life, call the \textbf{emergency number 112} first.
  	\item Enter the accident in the first-aid log (kept in the first-aid kit, located e.g.\ in the passage to the electronics workshop) and follow the further information given there.
  \end{itemize}
  \item Follow the instructions of the supervisors.
  \item \textbf{No open beverages} in the workshop rooms (risk of spilling, especially onto electronics and machines). Use only closable bottles or containers.
  \item \enquote{Legally binding} instructions are given only and exclusively by supervisors. A supervisor is someone who meets the requirements stated, among other places, in the \enquote{Mitmachen} (get involved) text and is listed as a supervisor in the list of supervisors on the website.
  \item Repairs are carried out only and exclusively by supervisors of the FAU FabLab or by companies or persons commissioned in writing by the FAU FabLab.
  \item If certain machines or tools are blocked for a long time, waiting visitors with shorter projects must be let in in between.
\end{itemize}
\vbox{\vspace{0,5cm}}

\section{Fire safety and emergencies}
\begin{itemize}
  \item In case of fire: stay calm, warn others, call the emergency number \textbf{112} and leave the building via the escape routes. Only if it is safe, fight an incipient fire with a fire extinguisher.
  \item On your first visit, find out where the fire extinguisher, escape routes and first-aid kit are. Ask a supervisor if you cannot find them.
  \item Escape routes and access to fire extinguishers and the first-aid kit must always be kept clear.
  \item Never leave hot tools (soldering iron, hot glue gun, etc.) unattended, and put them down only on designated non-flammable rests.
\end{itemize}

\vbox{\vspace{0,5cm}}

\section{Traffic light system}
Follow the traffic light labels on the devices!
\begin{itemize}
 \item RED: use only under the supervision of a supervisor (ask a supervisor)
 \item YELLOW: with a signed instruction
 \item GREEN: free to use
\end{itemize}

\vbox{\vspace{0,5cm}}

\section{Electronics}
\begin{itemize}
  \item If you want to use laboratory equipment (soldering station, power supply, etc.) in the electronics workshop, you must first sign an instruction on these rules
  \item \textbf{Do not use mains voltage!} Maximum 60\,V DC or 25\,V\subscript{rms} AC (extra-low voltage). This also applies to combinations of power supplies or batteries
 \item Do not eat, drink, apply makeup or similar while soldering. Wash your hands afterwards. (flux)
 \item Do not inhale solder fumes: use the extraction, if available, and ensure ventilation. The soldering iron is hot (risk of burns and fire): put it down only in its holder and switch it off after work.
 \item After use: switch off the switch at the table
 \item Put tools from the top drawers of the electronics workshop workstations back where they belong (color coding from tool to drawer)
\end{itemize}


\section{Metalworking}
\begin{itemize}
	\item Most tools are not suitable for stainless steel or hardened steel (e.g.\ polished shafts)
	\item When you are done:
	\begin{itemize}
		\item Sweep or vacuum up drilling chips
		\item Unclamp tools and put them away
		\item Put Proxxon tools back as well
	\end{itemize}
	\item Due to the lack of occupational safety equipment, fire protection and contamination in the FabLab, work such as welding, cutting and grinding with angle grinders is not possible in the FabLab (and of course not in the entire building either) and is therefore prohibited.
\end{itemize}


\subsection{Drill press, rotating tools}

\begin{itemize}
	\item Wear safety goggles when drilling brittle or splintering materials
	\item Never hold workpieces by hand. Change drill bits and tools only when the machine has stopped; for mains-powered machines, unplug them first.
	\item If necessary (small workpiece or large drill bit), secure the workpiece against turning
	\item Drill press: always use the machine vise and observe the notes posted there (exceptions: see the section for supervisors)
	\item \textbf{Secure long hair with a tight-fitting head covering (hairnet, cap or headscarf)!}
\end{itemize}
\textit{Reason: If hair is pulled into the drill, it can be torn out as a strand together with the hair root (the better case) or the entire scalp can be torn off (worst case).}

\begin{itemize}
	\item Keep away from anything that rotates (quickly)!
	\item \textbf{Wear neither gloves nor loose clothing} (e.g.\ smocks, aprons, scarves, ties, neckerchiefs, \dots) nor \textbf{jewelry} (rings, bracelets, watches, necklaces and the like).
\end{itemize}
\textit{Reason: If the machine catches clothing, it pulls body parts towards itself, with extremely unpleasant consequences.}

\begin{itemize}
	\item If in doubt, ask a supervisor
\end{itemize}



\section{Protective equipment (PPE) and dust}
\begin{itemize}
	\item The FabLab stores PPE for work (safety goggles, hearing protection); wear it when needed.
	\item Safety goggles are available in the marked containers at the workbench and at the etching table.
	\item Ear defenders hang at the milling machine.
	\item Avoid dust (wood, circuit boards): vacuum up dust instead of sweeping it. Always use power tools (jigsaw, sander) with a vacuum cleaner. The Festool vacuum cleaner (white-green) is approved for fine dust.
	\item Notify a supervisor if there is not enough protective equipment available.
	\item Put used items back after use.
\end{itemize}



\section{Paying for and using material from the FabLab}
\begin{itemize}
	\item Prices can be found at the checkout terminal or on the product labels
	\item Payment is made at the checkout terminal. If you have problems, contact a supervisor
\end{itemize}

\section{FabLab camera}
Photos taken with the FabLab camera are published online automatically. Before use, read the information sheet in the camera box.

\newpage
\section{Further notes for supervisors and key holders}
\textit{This section is intended only for supervisors and key holders.}

\subsection{Access and token}
Key authorizations (with FAUCard or token) are issued to \textbf{one} person and must not be passed on to other people.

\subsection{Leaving the FabLab}
If people without key authorization are present, you must hand over responsibility to another key holder who is present when leaving, or send everyone out of the FabLab.

When leaving, work through the checklist at the door (devices off, meeting room door closed, ...).

\subsection{Prohibited materials}
\begin{itemize}
	\item Asbestos (heat insulation in electrical appliances before 1993, e.g.\ toasters and hair dryers)
	\item Cathode ray tubes
	\item Self-built things connected to mains voltage without prior electrical inspection (tip: use a commercially available power supply instead)
\end{itemize}

\subsection{Fault reports}
In case of defects or unusual behavior: attach a fault notice to the device/tool (form \enquote{Function impaired} or \enquote{Out of order}), fill it in comprehensibly, and send an email to the mailing list.

\subsection{Drill press: exception to the machine vise}
Supervisors may deviate from the principle \enquote{use the machine vise} in individual cases, but only if it is ensured that nothing can happen. Where possible, do critical work with the cordless screwdriver.

\section{Copyright}

\textbf{Sharing and adapting this content is expressly encouraged}, you only have to (as a rule)
\begin{itemize}
 \item give the source and license correctly, for example:\\
FAU FabLab et al.: General Workshop Rules, \url{https://github.com/fau-fablab/werkstatt-einweisung},\\License CC-BY-SA 3.0, \url{https://creativecommons.org/licenses/by-sa/3.0/}.
 \item publish the new (entire) document under this license as well.
\end{itemize}

\footnotesize
\begin{wrapfigure}{r}{4cm}
\vspace{-1.5em}
\includegraphics[width=4cm]{fablab-document/cc-by-sa.pdf}
\vspace{-2.6em}
\end{wrapfigure}
This document \enquote{General Workshop Rules} by the FAU FabLab and other authors is, except for specially marked passages, licensed under a \emph{Creative Commons Attribution-ShareAlike 3.0 Unported} license. To view a copy of this license, visit \url{http://creativecommons.org/licenses/by-sa/3.0/}.

Source code and list of authors on GitHub: \url{https://github.com/fau-fablab/werkstatt-einweisung}.\\
Suggestions for improvement are welcome at \href{mailto:kontakt@fablab.fau.de}{kontakt@fablab.fau.de}.

Note on safety instructions: please note that occupational safety cannot be done by copy and paste; every instruction must be critically reviewed and adapted to its intended use.

This document comes from fau-fablab/werkstatt-einweisung@\Revision{}.
\normalsize

\end{document}
